\documentclass[11pt]{article}
\usepackage[letterpaper,margin=0.7in]{geometry}
\usepackage[T1]{fontenc}
\usepackage{lmodern}
\usepackage{amsmath,amssymb,tikz,array,fancyhdr}
\renewcommand\familydefault{\sfdefault}
\setlength\parindent{0pt}
\setlength\parskip{5pt}
\setlength{\headheight}{15pt}
\pagestyle{fancy}\fancyhf{}
\lhead{Week 2 / Lamp lab / Adult return-visit guide}
\fancyfoot[L]{Bellingham Math Circle / Week 2 / F02-RV-FAC-v1}
\fancyfoot[R]{\thepage}
\renewcommand\headrulewidth{0.4pt}
\newcommand\lead[1]{\textbf{#1}\par}
\begin{document}
\lead{Mathematical overview: what changes when the switches change}
A hidden switch toggles a fixed subset of lamps. Comparing its visible before and after states identifies that subset. Three individually observed presses identify three arbitrary switches. Three full-state experiments are necessary in the worst case even if several switches may be selected together: a three-lamp panel with three unknown masks has $2^9$ possibilities, while each observation has only $2^3$ outcomes. Fixed wiring and accurate observations are assumptions.

For a ring of $n\ge 3$ lamps, let switch $i$ toggle lamps $i,i+1,i+2$, wrapping around. If $3$ does not divide $n$, every target from all OFF has exactly one press subset. If $3$ divides $n$, only $2^{n-2}$ targets are reachable, and each has four press subsets. On six lamps, reachable targets have the same ON-count parity in $\{1,4\}$, $\{2,5\}$ and $\{3,6\}$. Repeated presses of one switch cancel in pairs.

An exchange switch moves the two endpoint counters into each other's positions; it does not turn them over. On a connected network it can make exactly the pictures with the same number of ON lamps. On a line, if the ON positions start at $a_1<\cdots<a_k$ and finish at $b_1<\cdots<b_k$, the exact minimum number of adjacent exchanges is $\sum_i|a_i-b_i|$. Disconnected networks require equal ON-counts in each component.

Experiments reveal candidates and counterexamples. Visible before/after differences explain wiring; a repeating press pattern explains the three-lamp obstruction; ordered particle displacements establish exchange distances. A sample of trials alone does not establish all-target reachability or impossibility.

\lead{Entries, preparation and flexible pacing}
The packet is shared Grades 2--5. Concrete observation and exchanging counters can suit K--1 with adult reading; finding hidden wiring without the adult deciding the trials requires keeping a before-state record. Grades 2--3 can conduct trials and seek routes. Grades 4--5 can organize complete cases and establish a lower bound. Counting to six is enough for entry; binary notation and linear algebra are adult descriptions, not child prerequisites.

Per pair: six reversible counters with clearly different faces, at most 20 mm across for the 21 mm working positions, pencil, one folder, and three switch cards for hidden wiring. Cut or tear the three-card strip from Page 1 in advance; one student-page copy can supply the cards while another stays intact as the board/record. Prepare five trays for the current eleven children: 30 counters, five folders, 15 cards; rotate the upper third child as operator/checker. Give each pair one page at a time. The hidden operator privately records the three fixed lamp subsets before any trial and changes every specified counter. For three-lamp switches the partner points to all three positions before turning counters; this prevents an imagined rule from disappearing into memory. Use paper position boards rather than real wiring.

Choose one investigation for 20--35 minutes, then return later. Keep three stable adult-led tables. Allow handling before the launch; adult reading and occasional recording help should leave children choosing experiments and routes. No obligation to finish the packet.

\lead{Launch together}
Let children handle the counters, agree which face is ON, and point out that a shaded printed lamp means ON and an empty one means OFF. Then demonstrate only the action for the chosen page. For hidden wiring, show an unrelated two-lamp switch and preserve the before picture while operating it. For triple switches, point to three consecutive lamps and turn all three over once. For transport, exchange two adjacent counters while keeping their faces up. Let a child make the next legal attempt. Keep the target visible; do not reveal the hidden three-switch panel or a complete ring solution.
\newpage
\lead{Problem 1: hidden wiring detective}
\textbf{Entry and timing.} Pairs alternate hidden operator and investigator, allowing 20--30 minutes for several panels. A young entry is noticing exactly which lamps changed after one press; predicting an untried combination is a later step.

\textbf{Setup.} The hidden partner chooses a fixed subset of lamps for each switch, including possibly none or all. Draw or mark these subsets on three cards kept behind a folder. The investigator sees all lamp states and chooses the switch tests. Keep one before/after drawing, or use a photograph/sketch, rather than concurrent tally systems. The operator may help point at labels but must not choose the investigator's tests.

\textbf{Exact solution.} Observe switch A alone, then switch B alone, then switch C alone, comparing each current state immediately before and after the press. A changed lamp belongs to that switch's subset; an unchanged lamp does not. Resetting is unnecessary if the before state is retained. After three such observations every subset is known, so any requested combined press can be predicted by toggling each lamp once per switch that includes it; two toggles cancel.

If cost means individual switch presses, three is necessary: a switch never pressed could have any wiring. If one experiment may press an arbitrary chosen set of switches once and then observe the whole panel, three is still necessary in the worst case by the $2^9$ versus $(2^3)^q$ count in the overview. Adaptively chosen tests do not evade that bound. A particular panel or extra prior knowledge can allow fewer tests; the bound is for guaranteed identification of arbitrary hidden wiring.

\textbf{Questions and hints.} ``Which lamps changed since this picture?'' If bookkeeping obstructs the idea, let the partner preserve a second before picture. Offer one-switch trials only after children have chosen some experiments. Ask whether a lamp that stayed OFF was disconnected from that switch or toggled an even number of times in a combined test.

\textbf{Continuations.} Hide four switches instead of three while keeping three lamps; guaranteed identification then requires four full-state trials. Restrict the possible masks and investigate whether fewer tests become enough. Introduce a faulty observation only as a separate later experiment with an explicit model; do not treat an operator mistake as mathematical evidence.

\textbf{Limits.} The panels have deterministic fixed toggling. No claim about noisy circuits, real electrical safety, or physical implementation is needed. Hidden wiring inference is new here; it does not repeat base target reachability or the earlier bridge-detective question.
\newpage
\lead{Problem 2: three lamps at a time}
\textbf{Entry and timing.} For every target, reset all counters to OFF and erase all press marks; then let pairs operate one numbered ring. A move changes the named lamp and the next two clockwise lamps, wrapping around. Allow 20--40 minutes and a later visit to compare ring sizes. Keep switch choices visible as a short move word or press markers.

\textbf{Exact classification.} Four lamps: all 16 targets are reachable, with one press subset each. Five lamps: all 32 targets are reachable, with one press subset each. Six lamps: only 16 of 64 targets are reachable, with four press subsets each. For six lamps the three paired groups $\{1,4\}$, $\{2,5\}$ and $\{3,6\}$ must have equal parities of their ON-counts. Exactly one ON lamp is impossible; all ON lamps are possible, in two presses: switches 1 and 4. The other all-ON press subsets are $\{2,5\}$, $\{3,6\}$ and all six switches. The nonempty do-nothing press subsets are $\{1,2,4,5\}$, $\{1,3,4,6\}$ and $\{2,3,5,6\}$. Symmetric differences with these give the four solutions of any reachable target.

\textbf{Why.} Let $x_i$ record whether switch $i$ was pressed an odd number of times. A lamp $i$ is toggled by $x_i,x_{i-1},x_{i-2}$. For a do-nothing subset, their sum is even. The recurrence forces $x_i$ to repeat every three positions: the possible patterns are 000, 011, 101, 110. Wrapping a ring whose length is not divisible by three forces the zero pattern, so the press-to-target map is one-to-one; finite equally sized input and output sets then make it onto. If the length is divisible by three there are four null patterns, so each reachable target has four preimages. On six lamps, every switch changes one member of each paired group, preserving equality of their three parities; the count shows that this condition is sufficient too.

\textbf{Printed targets.} On four lamps, ON only at 1 requires switches 1,3,4; all ON requires all four switches. On five lamps, ON only at 1 requires 2,3,5; all ON requires all five. On six lamps, ON only at 1 is impossible; all ON has the four solutions listed above, with minimum two presses.

\textbf{Questions and hints.} ``Which three lamps belong to this switch, including across the join?'' On six lamps suggest grouping opposite positions only after children encounter an impossible target. Let them invent a target they believe impossible and have the partner challenge it. A failed search alone is not an impossibility explanation.

\textbf{Return visits.} Make every target on a four-lamp ring, pooling the class records. Investigate nine lamps: there are 128 reachable targets among 512 and four press subsets per reachable target. Counting classes modulo three is a continuation; it is not necessary for the first concrete trials.

\textbf{Limits.} Exactly three consecutive distinct lamps change, including the selected starting lamp. Ring size is at least three. Repeating a switch is legal, but two repetitions cancel. This new rule changes the mathematical obstruction; ordinary pair-toggle parity must not be reused unchanged.
\newpage
\lead{Problem 3: exchanging the light}
\textbf{Entry and timing.} A partner physically exchanges the two adjacent counters; their faces stay unchanged. Reset to the same BEFORE picture for each line and ring trial and start a fresh move record. Begin with short lines and allow 20--35 minutes of finding shorter routes. Younger children can move a single ON counter toward a target; several counters and an optimality explanation require tracking order.

\textbf{Exact result and example.} The ON-count is conserved. On a line, match the first ON counter from the left to the first target ON position, the second to the second, and so on. Each must travel the absolute difference of its positions, giving a lower bound on exchanges. A route achieving the sum exists by moving counters in order without crossing. Exchanges of equal faces make no progress and are unnecessary in a shortest route.

For six lamps, ON at positions $\{1,4\}$ and target $\{3,6\}$ require $|1-3|+|4-6|=4$ exchanges. One shortest line route exchanges positions 4--5, 5--6, 1--2, 2--3. A target with three ON counters is unreachable from a two-ON start. These are adult demonstration/check examples; use the actual printed or child-created states for the main investigation.

\textbf{Printed rows (six lamps).} ON 1 to ON 6: line 5, ring 1. ON 1,2 to ON 5,6: line 8, ring 4. ON 1,3,5 to ON 2,4,6: line 3, ring 3. ON 1,2 to ON 2,4,6: impossible on both, since the ON-count changes. One four-swap ring route for the second pair is $1\!:\!6,\ 1\!:\!2,\ 5\!:\!6,\ 1\!:\!6$; notation $a:b$ exchanges those positions. A route of three swaps for the third pair is $1:2,3:4,5:6$. Line lower bounds use the ordered displacement formula. Ring minima were independently found by complete breadth-first search of each start's reachable states in the 64-picture model, so a shorter unseen route is excluded.

\textbf{Why all same-count states are possible.} On a line the ordered movement construction works. On a connected graph, a path of adjacent exchanges can transport a required ON counter to a missing target position; equivalently, edge swaps along a spanning tree can arrange any set of tokens. Repeatedly place required positions while keeping the remaining graph connected by removing leaves. The argument concerns existence, not an optimal graph distance formula. On a disconnected graph counts cannot move between components.

\textbf{Questions and hints.} ``Did the counter change face, or change place?'' Keep both the target and present state visible. If the pair tries to match crossing ON counters, ask which one is first from the left before and after each exchange. Do not demand a concurrent move tally: retain a move word and count later.

\textbf{Continuations.} Connect the ends into a ring and compare shortest routes: the ordered line formula cannot be applied without choosing the circular matching and direction correctly. Allow exchange along a branched network and distinguish reachability from shortest distance. Investigate indistinguishable counters versus individually labeled ones as a separate new condition.

\lead{Novelty and evidence}
The base Week 2 catalogs and existing bonus cover pair-toggle reachability/parity, line subsets, shortest or cheapest solutions, bridges, solo buttons and do-nothing sets. These three directions instead infer unknown wiring, change to triple switches, and transport existing states by exchanges. Sources and finite checks are supplied. The return-visit format is informed by \emph{Math Circle by the Bay}, Preface viii--x (PDF 9--11); these particular investigations are new designs, not copied problems. All are \textbf{unpiloted}; counter operation and staffing are \textbf{physically untested}. Record what children actually tried and which rule needed adult rescue before the next visit.

\end{document}
